\documentclass[10pt,a4paper]{amsart}
\usepackage[margin=19mm]{geometry}
\usepackage{amsmath,amssymb,mathtools}
\usepackage[scr=rsfs,cal=euler]{mathalfa}
\usepackage{xfrac}
\usepackage[unicode,hidelinks,bookmarks=false]{hyperref}
\newcommand{\ie}{i.e.\ }
\newcommand{\cf}{cf.\ }
\newcommand{\resp}{respectively\ }
\newcommand{\Subsection}{\subsection}
\providecommand{\nobreakdash}{\nobreak\hbox{-}}
\setlength{\emergencystretch}{2em}
\multlinegap0pt
\usepackage{amssymb}
\usepackage[shortlabels]{enumitem}
\usepackage{mathtools}
\usepackage{tikz}
\newtheorem{proposition}{Proposition}
\newtheorem{theorem}{Theorem}
\newcommand{\R}{{\mathbb R}}
\newcommand{\beq}{\begin{equation}}
\newcommand{\eeq}{\end{equation}}
\renewcommand{\ge}{\geqslant}
\renewcommand{\le}{\leqslant}
\title{Power law convergence and concavity for the Logarithmic Schrödinger equation}
\author{Marco Gallo, Sunra Mosconi, Marco Squassina}
\date{Source: arXiv:2411.01614v2 (CC BY 4.0)}
\begin{document}
\maketitle

\noindent\emph{Source and licence.} Power law convergence and concavity for the Logarithmic Schrödinger equation. Version arXiv:2411.01614v2 (CC BY 4.0), distributed by arXiv under the Creative Commons Attribution 4.0 International licence, http://creativecommons.org/licenses/by/4.0/, as recorded in the arXiv licence metadata for this exact version and shown on the abstract page https://arxiv.org/abs/2411.01614v2. Full licence notice: \texttt{licenses/arxiv-2411.01614-CC-BY-4.0.txt}.\par\smallskip

\noindent\textit{Editorial scope.} Counted unit: the article's theorem thm\_concav\_gs (source0.tex lines 1656--1661). The article's complete original proof of it is delivered with it (source0.tex lines 1663--1801, one \texttt{proof} environment of 139 lines). No mathematics is written, repaired or re-derived: the statement and the proof are the article's own lines, in the article's own notation, and no retained line is substituted or re-flowed.  Retained uncounted as the source-local context the proof needs: lines 261--286; lines 733--768; lines 1374--1389; lines 1434--1438.  Nothing was omitted from any retained span, so the package carries no omission record.  No retained line contains a citation, so the wrapper carries no bibliography and no bibliography entry.  No retained line contains a figure environment, an image-inclusion command or drawing code, so no image is delivered and the package declares no asset dependency.  Editorial formatting only, in addition: an AMS \texttt{amsart} preamble that restates the article's own declarations and macros used by the retained lines, namely the environments \texttt{proposition}, \texttt{theorem} on separate counters, together with the standard packages those lines need.  The retained environments print in this document as Proposition 1, Theorem 1 in order of appearance, whereas the article prints the same items with the numbering of its own full document.  The complete native source of the article as distributed in the arXiv e-print for this exact version is bundled unchanged as \texttt{sources/167892/source0.tex}.
The main goal of this paper is to study the following \emph{nondispersive Logarithmic Schrödinger} equation
\beq
\label{LSeq}
\begin{cases}
-\Delta u=u\, \log u^2& \text{in $\Omega$}\\
u>0&\text{in $\Omega$}\\
u=0&\text{on $\partial\Omega$}
\end{cases}
\eeq
on some $\Omega \subset \R^N$ bounded, $N\geq 1$, together with the following \emph{Lane-Emden} equation
\beq
\label{LEeq}
\begin{cases}
-\Delta u=\sigma\, (u^q-u)&\text{in $\Omega$}\\
u>0&\text{in $\Omega$}\\
u=0&\text{on $\partial\Omega$}
\end{cases}
\eeq
when $\sigma>0$ and $q \in\ ]1, 2^*-1[$, $2^*:=2N/(N-2)$ if $N\geq 3$ or $2^*:=\infty$ if $N\le 2$.
In particular, following the heuristic 
\beq
\label{heu}
-\Delta u=\frac{2}{q-1} (u^q-u) \quad \longrightarrow\quad -\Delta u= u\log u^2
\eeq
 as $q \to 1^+$, we aim to study the convergence of the solutions of \eqref{LEeq} with $\sigma = 2/(q-1)$ to a solution of \eqref{LSeq}. 
When the domain $\Omega$ is convex, we investigate the concavity properties of solutions to \eqref{LEeq} and, as a byproduct of the aforementioned convergence, we obtain $\log$-concavity for a solution of \eqref{LSeq}, which is our main result.

To state it we recall some elementary facts on the second fundamental form of level sets of a $v\in C^2(A)$, where $A$ is an open subset of $\R^N$ and $N\ge 2$. 
Let $t \in v(A)$ and $x_0 \in \{v=t\}$; in order for the level set $\{v=t\}$ to be a $C^2$ submanifold near $x_0$ we will suppose that $Dv(x_0)\ne 0$. 
We choose a normal vector to $\{v=t\}$ at $x_0$ as
\beq
\label{normale}
 {\bf n}_{x_0} :=-\frac{Dv(x_0)}{|Dv(x_0)|}
 \eeq
 and set the tangent space $T_{x_0}:=\{z\in \R^N:\langle {\bf n}_{x_0}, z\rangle=0\}$. 
The \emph{second fundamental form} of the level set of $v$ at $x_0$ is the quadratic form defined on $T_{x_0}$ as
\[
{\rm II}_{x_0}(v)(z):=\langle {\bf n}_{x_0}, \ddot{\gamma}(0)\rangle
\]
where $\gamma:\ ]-\delta, \delta[\ \to \{v=t\} $ is a curve (for some $\delta>0$) parametrised by arc-length such that $\gamma(0)=x_0$ and $\dot\gamma(0)=z$. 
This definition is independent of the choice of $\gamma$, under the required constraints. 
We can differentiate twice the relation $v(\gamma(s))=v(x_0)$ to get
\[
\langle D^2v(\gamma(s))\,\dot\gamma(s), \dot\gamma(s)\rangle +\langle Dv(\gamma(s)), \ddot{\gamma}(s)\rangle=0
\]
for all sufficiently small $|s|$. 
Recalling \eqref{normale} (so that $\partial_n v(x_0)=-|Dv(x_0)|$), it follows that
\beq
\label{secondff}
{\rm II}_{x_0}(v)(z)=\frac{\langle D^2v(x_0)\,z, z\rangle}{|Dv(x_0)|}.
\eeq
The \emph{mean curvature} $H_{x_0}(v)$ of $\{v=t\}$ at $x_0$ is the trace of the second fundamental form. 
As such, it is the sum of the {\em principal curvatures} of $\{v=t\}$ at $x_0$, each of the latter being computed as ${\rm II}_{x_0}(v)(z_i)$ where $\{z_i\}_{i=1}^{N-1}$ is a family of $N-1$ orthonormal eigenvectors in $T_{x_0}$ for $D^2{\rm II}_{x_0}(v)$. 
Hence 
\[
H_{x_0}(v):=\sum_{i=1}^{N-1} {\rm II}_{x_0}(v)(z_i)=\frac{1}{|Dv(x_0)|}\sum_{i=1}^{N-1}\langle D^2v(x_0)\,z_i, z_i\rangle;
\]
we complete such system by setting $z_N:={\bf n}_{x_0}$, to obtain an orthonormal basis of $\R^N$. 
Since the Laplacian is invariant under unitary change of variables, we obtain
\beq
\label{curvmedlap}
\Delta v(x_0)=\sum_{i=1}^N\langle D^2v(x_0)\,z_i, z_i\rangle=\frac{\langle D^2v(x_0)\,Dv(x_0), Dv(x_0)\rangle}{|Dv(x_0)|^2}+H_{x_0}(v)\, |Dv(x_0)|.
\eeq

\begin{proposition}[Convergence to eigenfunction]
\label{prop_conv_eigenf}
Let $\Omega$ be bounded and convex, $\sigma>0$, and $u_n$ a solution of \eqref{LEeq} for a sequence $q_n\to 1^+$. 
Denote by $\varphi_1$ the first positive eigenfunction of the Dirichlet Laplacian normalised so that $\|\varphi_1\|_\infty=1$ and by $\lambda_1$ the corresponding first eigenvalue. 
Then 
\begin{align}
&\label{limv1}
\frac{u_{n}}{\|u_n\|_\infty} \to \varphi_1\quad \text{ in $C^2_{\rm loc}(\Omega)\cap C^0(\overline{\Omega})$},
\\&\label{limMq}
 \|u_n\|_\infty^{q_n-1} \to 1+\frac{\lambda_1}{\sigma},
\\& \label{limuq-1}
u_n^{q_n-1}\to 1+\frac{\lambda_1}{\sigma}\quad \text{ in $C^0_{{\rm loc}}
(\Omega)$}.
\end{align}
Moreover, if $\partial\Omega$ is smooth, then the convergence in \eqref{limv1} holds in $C^{2}(\overline{\Omega})$ as well.
\end{proposition}

\begin{proposition}[Uniqueness]
\label{corol_uniq}
Let $\Omega$ be bounded and convex and $\sigma>0$. 
Then there exists $q_0=q_0(\sigma, \Omega)>1$ such that for any $q\in \, ]1, q_0]$, \eqref{LEeq} has a unique solution.
\end{proposition}

\begin{theorem}[Concavity near $q=1$]
\label{thm_concav_gs}
Let $\sigma>0$ and $\Omega\subseteq \R^N$ be bounded, smooth and strongly convex if $N\ge 2$. 
Then there exists $q_0=q_0(\sigma, \Omega) \in\ ]1, 2^*-1[$ such that, for $q \in\ ]1, q_0[$ 
the solution $u_{q, \sigma}$ to \eqref{LEeq} is unique (indeed, it is a ground state $u_{q, \sigma} \in \mathcal{GS}_{q, \sigma}$), strongly $\log$-concave, and thus strongly $(1-q)/2$-concave in $\Omega$.
\end{theorem}

\begin{proof}
Uniqueness has been proven in Proposition \ref{corol_uniq} for $q-1$ sufficiently small, which we will assume henceforth. 
Note that for $q\in\ ]1, 2^*-1[$ any solution $u_{q, \sigma}$ of \eqref{LEeq} produces, by considering $u_{q, \sigma}(\sqrt{\sigma}\, \cdot)$, a solution of \eqref{LEeq} for $\sigma=1$ on the domain $\Omega/\sqrt{\sigma}$. 
Being $\sigma$ fixed, we can suppose that $\sigma=1$ and omit henceforth the dependence on $\sigma$. 
 
We set $v_q:=u_q/\|u_q\|_\infty$ and note that 
 \[
 D^2\log u_q= D^2\log v_q=\frac{ 1}{v_q}\,\left[ D^2v_q-\frac{1}{v_q}\, Dv_q\otimes Dv_q\right].
 \]

 Thus we focus on the matrix 
\begin{equation}\label{eq_def_Mq}
 M(v_q):= 
 \frac{1}{v_q}\, Dv_q\otimes Dv_q -D^2v_q,
\end{equation}
which we will prove to be positive definite.

\smallskip
 {\em Step 1: bound in the normal directions}.\ \\
For sufficiently small $\delta>0$ let $\Phi_t:[0, \delta[\,\times \partial\Omega\to \Omega$ be a $C^1$ (in both $t$ and $x$) family of diffeomorphisms from $\partial\Omega$ to $\{{\rm dist}(x, \partial\Omega)=t\}$, and let $n$ denote the corresponding $C^1$ extension of the interior normal to $\partial\Omega$, defined on 
\[
\Omega_\delta :=\{x\in \Omega:{\rm dist}(x, \partial\Omega)<\delta\}.
\]
 Note that any $\xi\in \R^N$ such that $(n(x), \xi)=0$ is the image of a unique tangent vector $\xi'$ to $\partial\Omega$ at the point $\Phi^{-1}_{{\rm dist}(x, \partial\Omega)}(x)$ and the corresponding map is $C^1$. 

Let $\varphi_1$ be the first positive eigenfunction of the Dirichlet Laplacian such that $\|\varphi_1\|_\infty=1$. 
 The Hopf Lemma ensures that there exist $\delta, \theta>0$ such that 
 \[
 \inf_{\Omega_{\delta}}\frac{\partial \varphi_1}{\partial n}\ge 3\, \theta
 \]
 where $n$ is the interior normal to $\partial\Omega$.
Let $q_0\in \ ]1, 2^*-1[$, given in Proposition \ref{corol_uniq}, be such that, for any $1<q\le q_0$, \eqref{LEeq} has a unique solution $u_q$ (which is a ground state).
 By the $C^1(\overline{\Omega})$ convergence $v_q\to \varphi_1$ as $q\to 1^+$ proved in Proposition \ref{prop_conv_eigenf}, there exists $1<q_0'\le q_0 $ such that 
 \beq
 \label{norm}
 \inf_{\Omega_\delta}\frac{\partial v_q}{\partial n}\ge 2\, \theta \qquad \text{for all $q\in \ ]1, q_0']$}.
 \eeq
 Since, again by Proposition \ref{prop_conv_eigenf}, there exists $C>0$ such that
 \begin{equation}\label{eq_unif_c2_est}
 \|v_q\|_{C^2(\overline{\Omega})} \leq C \qquad \text{for all $q\in \ ]1, q_0']$},
 \end{equation}
inequality \eqref{norm} gives
 \beq
 \label{norm2}
 \left(M(v_q)\, n, n\right)\ge \frac{1}{v_q}\, \left(\frac{\partial v_q}{\partial n}\right)^2-|D^2 v_q|\ge \frac{4\, \theta^2}{v_q} -C\ge \frac{4\, \theta^2}{C\, \delta}- C
 \eeq
in $\Omega_\delta$. 
This concludes the proof when $N=1$, so we will suppose from now on that $N\ge 2$.

\smallskip
 {\em Step 2: bound in the tangential directions}.\ \\
 Let ${\rm II}_{x}(\partial\Omega)$ be the second fundamental form of $\partial\Omega$ with respect to the inner normal direction $ n_x$ at a point $x\in \partial\Omega$, so that by assumption there exists $k_0>0$ such that 
\beq
\label{strc}
{\rm II}_{x}(\partial\Omega)(\xi)\ge k_0\, |\xi|^2, \qquad \forall \, x\in \partial\Omega, \, \xi\bot n_x.
\eeq
Since $\partial\Omega=\{-v_q=0\}$ and $-v_q<0$ in $\Omega$, it holds
\[
{\rm II}_{x}(\partial\Omega)={\rm II}_{x}(-v_q)=-{\rm II}_{x}(v_q)
\]
for all $x\in \partial\Omega$.
From the latter, \eqref{secondff}, \eqref{strc} and \eqref{norm} we have
\[
\left(D^2v_q(x)\, \xi, \xi\right)=|Dv_q(x)|\, {\rm II}_x(v_q)(\xi)=- \frac{\partial v_q(x)}{\partial n}{\rm II}_{x}(\partial\Omega)(\xi) \le -2\, k_0\, \theta\, |\xi|^2
\]
for all $x\in \partial\Omega$, $\xi\bot n_x$ and $q\in \ ]1, q_0']$.

 Since $Dv_q\otimes Dv_q\ge 0$, $n\in C^1(\Omega_\delta)$ and $v_q$ is uniformly bounded in $C^2(\overline{\Omega})$ we infer that for any sufficiently small $\delta$ it holds
 \beq
 \label{tang}
 \left(M(v_q)(x)\, \xi, \xi\right)\ge k_0\, \theta\, |\xi|^2
 \eeq
 for any $q\in \ ]1, q_0']$, $x\in \Omega_\delta$ and $\xi=\xi(x)$ such that $\xi(x)\bot n(x)$. 
 
 \smallskip
 {\em Step 3: bound in the mixed directions}.\ \\
 Finally, since $(Dv_q(x_0), \xi)=0$ for all $x_0\in \partial\Omega$ and $\xi\bot n_{x_0}$ with $|\xi|=1$ and the boundedness of $v_q$ in $C^2(\overline{\Omega})$ for all $q\in \ ]1, q_0']$, we deduce through the Lipschitz character of $\xi\mapsto \xi'\in T_{\partial \Omega}$ the uniform bound
 \[
|(Dv_q(x), \xi(x))|\le C\, {\rm dist}(x, \partial\Omega)
 \]
for all $x\in \Omega_\delta$ and $\xi\bot n(x)$ with $|\xi|=1$. In particular, since by \eqref{norm} it holds
\[
{\rm dist}(x, \partial\Omega)\le C\, v_q(x)
\]
in $\Omega_\delta$, for a constant $C$ independent of $q\in \ ]1, q_0']$, we get 
\begin{equation}\label{eq_mixed_dir}
\left(M(v_q)\, \xi, n\right)\le |D^2 v_q|\, |\xi|+ \frac{1}{v_q} \left|\frac{\partial v_q}{\partial n}\right|\, \left|(Dv_q, \xi)\right|\le C\, |\xi|
\end{equation}
for all $x\in \Omega_\delta$, $\xi\bot n(x)$ and $q\in \ ]1, q_0']$.

 \smallskip
 {\em Step 4: convexity in $\Omega_{\delta}$}.\ \\
Writing any vector as $\xi+t\, n$ for $\xi\bot n$ and $t\in \R$, it follows from \eqref{norm2}, \eqref{tang} and \eqref{eq_mixed_dir}, that in $\Omega_\delta$ we have
 \[
 \begin{split}
 \left(M(v_q) \, (\xi+tn), \xi+tn\right)&=(M(v_q)\, \xi, \xi) +2\, t\, (M(v_q)\, \xi, n)+t^2\, (M(v_q) \, n, n)\\
 &\ge 
 k_0\, \theta\, |\xi|^2-C\, t\, |\xi| +t^2\left(\frac{4\, \theta^2}{C\, \delta}- C\right)
 \end{split}
 \]
 for all $q\in \ ]1, q_0']$. It suffices to choose $\delta>0$ sufficiently small (depending only on the parameters and thus not on $q$) to obtain a positive constant $\theta_0'$ such that
 \[
\left( M(v_q)\, z, z\right)\ge \theta_0'\, |z|^2 \quad \hbox{in $\Omega_{\delta}$},
\]
for all $q\in \ ]1, q_0']$. Then since $v_q\le C\, \delta $ in $\Omega_\delta$, we find
 \[
 D^2\log v_q=-\frac{M(v_q)}{v_q}\le -\frac{\theta_0'}{C\, \delta}\, {\rm Id} 
 \]
in $\Omega_\delta$, for all $q\in \ ]1, q_0']$. 
 
 \smallskip
 {\em Step 5: conclusion}.\ \\
Recall that $\log v_q\to \log \varphi_1$ in $C^2(\Omega\setminus\Omega_\delta)$ by Proposition \ref{prop_conv_eigenf}. Note that $D^2\log\varphi_1$ is negative definite everywhere in $\Omega$ by a classical application of the constant rank theorem, hence $\log\varphi_1$ is locally strongly concave in $\Omega$. 
By $C^2$-convergence, this ensures that for a sufficiently small $q_0''>1$ and $\theta_0''>0$ it holds
 \[
 D^2\log v_q\le -\theta_0''\, {\rm Id}
 \]
 in $\Omega\setminus \Omega_\delta$ for all $q\in \ ]1, q_0'']$. 
All in all we have proved that for a constant $\theta_0=\min\{\theta_0', \theta_0''\}>0$ and $q_0=\min\{q_0', q_0''\}>1$, the inequality
 \beq
 \label{d2l}
 D^2\log u_q=D^2\log v_q\le -\theta_0\, {\rm Id}
 \eeq
 holds true in $\Omega$ for all $q\in \ ]1, q_0]$.

 To prove the final assertion, we compute 
 \[
 \begin{split}
 D^2 u_q^{\frac{1-q}{2}}&=\frac{q-1}{2\, u_q^{\frac{q+1}{2}}}\left(\frac{q+1}{2}\frac{D u_q\otimes D u_q}{u_q}- D^2u_q\right)\\
 &=\frac{q-1}{2\, u_q^{\frac{q+1}{2}}}\left(\frac{q-1}{2}\frac{D u_q\otimes D u_q}{u_q}-u_q\, D^2\log u_q\right)
 \end{split}
 \]
 and note again that the first matrix term is non-negative definite. 
Thus from \eqref{d2l} we have
 \[
 D^2 u_q^{\frac{1-q}{2}}\ge -\frac{q-1}{2}\, u_q^{\frac{1-q}{2}} D^2\log u_q\ge \theta_0\, \frac{q-1}{2}\, \|u_q\|_\infty^{\frac{1-q}{2}}\, {\rm Id}
 \]
and $u_q^{(1-q)/2}$ is strongly convex on $\Omega$ for all $q\in \ ]1, q_0]$. 
\end{proof}

\end{document}
